\section{Points, idempotents and Green's relations}\label{sec:points}

Deck: right $M$-sets as presheaves on $\B M$. Arch: Bridge~\ref{br:pts}, essential points are the objects of $\Kar(\B M)^{\op}$, which by Example~\ref{ex:karM} are the idempotents of $M$ up to Green's relation $\mathcal D$.

\subsection{The theorem}

\begin{theorem}\label{th:points}
Let $M$ be a monoid. For an idempotent $e = e^2$ of $M$ and a right $M$-set $X$, put $Xe = \{x\in X : x\cdot e = x\}$; a map of $M$-sets $X\to Y$ sends $Xe$ into $Ye$, so $X\mapsto Xe$ is a functor to $\Set$.
\begin{enumerate}
\item The functors from right $M$-sets to sets that preserve all limits and all colimits are, up to isomorphism, exactly the functors $X\mapsto Xe$.
\item $X\mapsto Xe$ and $X\mapsto Xf$ are isomorphic iff there are $a\in fMe$ and $b\in eMf$ with $ba = e$ and $ab = f$ (i.e.\ $e\mathrel{\mathcal D}f$).
\item The natural transformations from $X\mapsto Xf$ to $X\mapsto Xe$ correspond to the elements of $fMe$, an element $a$ acting by $x\mapsto x\cdot a$.
\end{enumerate}
\end{theorem}

\inwords{} The \q{ways of looking at all $M$-sets at once} that respect all constructions are given by the idempotents: look at the elements fixed by $e$. Two idempotents give the same way of looking iff they are related by Green's relation $\mathcal D$. In Galois-theoretic language, these are the \q{fibre functors} of the category of $M$-sets.

\gen{} Bridge~\ref{br:pts}, together with the facts that a functor $\PSh(\B M)\to\Set$ that preserves all colimits has a right adjoint and one that preserves all limits (and colimits) has a left adjoint, so that such functors are exactly the essential points. The object $e$ of $\Kar(\B M)$ corresponds to the retract $eM$ of the representable presheaf $M$, and to the functor $\Hom(eM,-)$; and $\Hom(eM,X)\cong Xe$ by $\varphi\mapsto\varphi(e)$.

\direct{} We check the parts that do not need the adjoint functor theorem. The functor $X\mapsto Xe$ preserves limits, since $Xe\cong\Hom(eM,X)$, and colimits: a colimit of $M$-sets is computed on underlying sets, as a quotient of a disjoint union, and applying $x\mapsto x\cdot e$ to representatives and to the generating identifications shows that every $e$-fixed class contains an $e$-fixed representative ($x\cdot e$) and that two $e$-fixed elements are identified by a chain of $e$-fixed elements. For (3): a natural transformation $\Hom(fM,-)\to\Hom(eM,-)$ is, by the Yoneda lemma, a map $eM\to fM$, i.e.\ the element $a = $ image of $e$, which satisfies $a = a e$ and $a\in fM$, so $a\in fMe$; it acts on $x\in Xf$ by $x\mapsto x\cdot a$. For (2): an isomorphism is a pair $a\in fMe$, $b\in eMf$ whose composites are the identities, $ba = e$ and $ab = f$. Part (1), that there are no other such functors, is the content of Bridge~\ref{br:pts}: a colimit-preserving functor $p$ is determined by $p(M)$ with its left $M$-action, a limit-preserving one is representable by some $P$, and $P$ must be a retract of $M$, of the form $eM$.\qed

\status{} Essential points of $\PSh(\B M)$ correspond to idempotents \cite[Cor.~4.2]{rogers-1905}; isomorphism of objects of the Karoubi envelope of a semigroup is the relation $\mathcal D$ \cite[Prop.~3.1, \S3.2]{costa-steinberg}. We have not found the combination stated in one place, nor the counts below.

\subsection{Counting the points}

\begin{corollary}\label{co:countpoints}
Up to isomorphism, the functors from right $M$-sets to sets preserving all limits and colimits are:
\begin{enumerate}
\item for a group $G$: exactly one, the forgetful functor $U$, and $\Aut(U)\cong G$;
\item for the multiplicative monoid $M_n(k)$ of $n\times n$ matrices over a field: exactly $n+1$, one for each rank $r = 0,\dots,n$;
\item for the monoid $T_n$ of all maps of an $n$-element set: exactly $n$, one for each image size $1,\dots,n$;
\item for the bicyclic monoid: exactly one;
\item for a commutative monoid: one for each idempotent, pairwise non-isomorphic; for instance, for a finite lattice under $\wedge$, one for each element.
\end{enumerate}
\end{corollary}

\begin{proof}
(1) The only idempotent of $G$ is $1$, and $X\cdot 1 = X$. By Theorem~\ref{th:points}(3) the endomorphisms of $U$ are the maps $x\mapsto x\cdot g$, $g\in G$, and $x\mapsto x\cdot g$ followed by $x\mapsto x\cdot h$ is $x\mapsto x\cdot gh$; so $g\mapsto(x\mapsto x\cdot g^{-1})$ is an isomorphism $G\cong\Aut(U)$.
(2) Example~\ref{ex:karMat}: projections are $\mathcal D$-related iff they have the same rank, and there are projections of each rank $0,\dots,n$.
(3) Exercise~\ref{exo:karTn}; idempotents of $T_n$ exist with every image size from $1$ to $n$.
(4) The idempotents of the bicyclic monoid are $e_i = q^ip^i$. The elements $a = q^i\in e_iMe_0$ and $b = p^i\in e_0Me_i$ satisfy $ba = p^iq^i = 1 = e_0$ and $ab = e_i$. So all idempotents are $\mathcal D$-related to $1$.
(5) In a commutative monoid, $ba = e$ and $ab = f$ force $e = f$.
\end{proof}

\begin{example}[What the $n+1$ points of $M_n(k)$-sets see]
Let $M = M_n(k)$ act on the row vectors $X = k^n$ by right multiplication, $x\cdot m = xm$. For the projection $e_r = \operatorname{diag}(1,\dots,1,0,\dots,0)$ of rank $r$, $Xe_r = \{x : xe_r = x\}$ is the set of row vectors supported on the first $r$ coordinates, a copy of $k^r$. The $n+1$ points of the topos of $M_n(k)$-sets look at this $M$-set through the $n+1$ \q{windows} $k^0, k^1,\dots,k^n$.
\end{example}

\begin{remark}
Corollary~\ref{co:countpoints}(1) is the reconstruction of a group from its category of $G$-sets together with the forgetful functor, the simplest case of Grothendieck's Galois theory. Its generalisation (2)--(5) replaces \q{the} fibre functor by a family of fibre functors indexed by $\mathcal D$-classes of idempotents.
\end{remark}

\begin{exercise}\label{exo:Nmult}
Let $M = (\mathbb N,\cdot)$, the natural numbers including $0$ under multiplication. Find its idempotents and the essential points of the topos of right $M$-sets. Is this topos Boolean? De Morgan?
\end{exercise}
